% SolarMax — Testing the solar panel with the AstroAI AM33D multimeter
% Compile:  pdflatex panel_test_am33d.tex   (run twice for the table of contents)

\documentclass[11pt]{article}
\usepackage[margin=2.2cm]{geometry}
\usepackage{xcolor}
\usepackage[most]{tcolorbox}
\usepackage{enumitem}
\usepackage{titlesec}
\usepackage{listings}
\usepackage{array}
\usepackage[colorlinks=true,linkcolor=blue!60!black,urlcolor=blue!60!black]{hyperref}

% ── Colors / section styling ─────────────────────────────────────────────────
\definecolor{brandblue}{HTML}{1565C0}
\definecolor{brandgreen}{HTML}{2E7D32}
\titleformat{\section}{\Large\bfseries\color{brandblue}}{\thesection.}{0.5em}{}

% ── Callout boxes ─────────────────────────────────────────────────────────────
\newtcolorbox{tip}{colback=green!6,colframe=brandgreen,boxrule=0.8pt,
  left=8pt,right=8pt,top=6pt,bottom=6pt,title=\textbf{Tip},fonttitle=\color{white}}
\newtcolorbox{warn}{colback=red!5,colframe=red!70!black,boxrule=0.8pt,
  left=8pt,right=8pt,top=6pt,bottom=6pt,title=\textbf{Important},fonttitle=\color{white}}
\newtcolorbox{note}{colback=blue!5,colframe=brandblue,boxrule=0.8pt,
  left=8pt,right=8pt,top=6pt,bottom=6pt,title=\textbf{Note},fonttitle=\color{white}}

\lstdefinestyle{cmd}{basicstyle=\ttfamily\small,backgroundcolor=\color{black!6},
  frame=single,rulecolor=\color{black!25},framesep=6pt,aboveskip=6pt,belowskip=6pt,
  xleftmargin=4pt,xrightmargin=4pt,breaklines=true,columns=fullflexible,keepspaces=true}
\lstset{style=cmd}

\setlist[enumerate]{leftmargin=1.7em,itemsep=3pt,topsep=4pt}
\renewcommand{\arraystretch}{1.3}

\begin{document}

% ── Title ─────────────────────────────────────────────────────────────────────
\begin{center}
{\Huge\bfseries\color{brandblue} SolarMax}\\[4pt]
{\LARGE Testing the Solar Panel --- AM33D Multimeter Only}\\[8pt]
{\large Measure the panel tracking vs.\ not tracking. No resistors, no power supply.}\\[6pt]
\rule{0.9\linewidth}{0.8pt}
\end{center}

\vspace{4pt}
\tableofcontents
\vspace{8pt}

% =============================================================================
\section{The idea}
You can evaluate the panel with \textbf{just the AM33D multimeter} by taking two
readings on the bare panel:
\begin{itemize}[leftmargin=1.6em]
  \item \textbf{Voc} --- the open-circuit \textbf{voltage} (nothing connected but the meter).
  \item \textbf{Isc} --- the short-circuit \textbf{current} (meter straight across the panel).
\end{itemize}

\begin{note}
\textbf{Why this shows the tracker working:} a panel's \textbf{short-circuit current
(Isc) is proportional to how much sunlight lands on it} --- so it rises when the panel
faces the sun squarely and falls when it doesn't. Voltage (Voc) barely changes with
aim. \textbf{So Isc is the number that proves tracking helps.}
\end{note}

The panel is a Renogy 50\,W (RNG-50D-SS): in full sun expect \textbf{Voc\,$\approx$\,22\,V}
and \textbf{Isc up to $\approx$\,3\,A}. Rough power: $P \approx 0.74 \times V_{oc} \times I_{sc}$.

% =============================================================================
\section{Safety \& setup (read first)}
\begin{itemize}[leftmargin=1.6em]
  \item Do this outside in \textbf{steady, strong sun} (clouds change the reading).
  \item \textbf{The panel does not connect to your power supply} --- it makes its own
        power from light. Only the panel and the meter are involved.
\end{itemize}

\begin{warn}
\textbf{Shorting a solar panel is safe} (unlike a battery) --- it can only push its
\emph{Isc}, about 3\,A. But two AM33D rules matter:
\begin{itemize}[leftmargin=1.4em,topsep=2pt]
  \item For current, use the \textbf{10\,A jack + 10\,A dial position}. Do \textbf{not}
        use the \texttt{V$\Omega$mA} jack for this --- it's fused at 250\,mA and 3\,A
        would blow that fuse.
  \item The datasheet limits the 10\,A range to \textbf{15 seconds at a time}. Take the
        Isc reading quickly, then remove the probes.
\end{itemize}
\end{warn}

% =============================================================================
\section{AM33D quick settings}
Black lead is always in \textbf{COM}. You'll use exactly two setups:

\begin{center}
\begin{tabular}{|>{\bfseries}l|l|l|l|}
\hline
Reading & Dial position & Red lead jack & Typical display \\ \hline
Voc (voltage) & \textbf{200} in the DC-V block (V with $=$) & \texttt{V$\Omega$mA} & \texttt{21.8} (volts) \\ \hline
Isc (current) & \textbf{10A} in the DC-A block & \texttt{10A} & \texttt{2.85} (amps) \\ \hline
\end{tabular}
\end{center}

\begin{note}
Use the \textbf{DC} side (the ``V'' with a solid line above a dashed line, \(V\!=\)),
\emph{not} the AC ``V\raisebox{0.2ex}{$\sim$}''. If the voltage shows a leading
\textbf{``$-$''}, your leads are just reversed --- no harm, the number is still correct.
\texttt{OL} means over-range (shouldn't happen on these ranges here).
\end{note}

% =============================================================================
\section{Measurement 1 --- Open-circuit voltage (Voc)}
\begin{enumerate}
  \item Red lead in the \texttt{V$\Omega$mA} jack, black lead in \texttt{COM}.
  \item Turn the dial to \textbf{200} in the \textbf{DC voltage} block.
  \item Touch the red probe to the panel's \textbf{+} wire and black to \textbf{$-$}.
  \item Read the voltage. In full sun it should be around \textbf{18--22\,V}.
\end{enumerate}

\begin{tip}
Press \textbf{HOLD} to freeze a reading so you can write it down. Voc changes only a
little with aim --- that's expected; the action is in the current (next).
\end{tip}

% =============================================================================
\section{Measurement 2 --- Short-circuit current (Isc)}
\begin{enumerate}
  \item Move the red lead to the \textbf{10\,A} jack (black stays in \texttt{COM}).
  \item Turn the dial to \textbf{10A} in the \textbf{DC current} block.
  \item Touch the probes across the panel's \textbf{+} and \textbf{$-$} wires --- this
        briefly shorts the panel \emph{through the meter}, which is fine.
  \item Read the current \textbf{within a few seconds}, then lift the probes. Full sun,
        aimed at the sun: up to about \textbf{3\,A}.
\end{enumerate}

\begin{warn}
Keep each Isc reading \textbf{under 15 seconds}. When you're done measuring current,
\textbf{move the red lead back to the \texttt{V$\Omega$mA} jack} so you don't
accidentally short something later while the meter is in high-current mode.
\end{warn}

% =============================================================================
\section{The experiment --- tracking vs.\ not tracking}
Do both readings twice, back-to-back, in the same sunlight:

\begin{enumerate}
  \item \textbf{Tracking:} panel aimed \textbf{straight at the sun} (tracker running and
        pointed at it, or aim it by hand so its face is square to the sun). Record Voc,
        then Isc.
  \item \textbf{Not tracking:} leave the panel at a \textbf{fixed, off-sun angle} (e.g.\
        flat/horizontal, or rotated $\sim$45$^\circ$ away). Record Voc, then Isc.
  \item Repeat the pair at a few times of day (morning, noon, afternoon) --- the gap is
        biggest when the sun is low.
\end{enumerate}

\begin{center}
\begin{tabular}{|l|c|c|c|c|c|}
\hline
\textbf{Time} & \multicolumn{2}{c|}{\textbf{Tracking}} & \multicolumn{2}{c|}{\textbf{Not tracking}} & \textbf{Isc gain} \\ \cline{2-5}
 & Voc (V) & Isc (A) & Voc (V) & Isc (A) & (\%) \\ \hline
Morning  & \quad\quad & \quad\quad & \quad\quad & \quad\quad & \quad\quad \\ \hline
Noon     & & & & & \\ \hline
Afternoon& & & & & \\ \hline
\end{tabular}
\end{center}

% =============================================================================
\section{What the numbers mean}
\begin{itemize}[leftmargin=1.6em]
  \item \textbf{Isc gain} $=\dfrac{I_{sc}(\text{tracking}) - I_{sc}(\text{not})}{I_{sc}(\text{not})}\times 100\%$.
        A positive gain is the tracker capturing more sunlight --- that's your result.
  \item \textbf{Approx.\ power} for any reading: $P \approx 0.74 \times V_{oc}\times I_{sc}$.
        Compare the tracking vs.\ not-tracking power at each time.
  \item You should see \textbf{Voc nearly unchanged} but \textbf{Isc (and power) clearly
        higher when tracking} --- especially morning and afternoon.
\end{itemize}

\begin{note}
For a fair comparison, take each tracking / not-tracking pair \textbf{quickly and in the
same light}. If a cloud passes, redo that pair.
\end{note}

% =============================================================================
\section{Troubleshooting}
\begin{itemize}[leftmargin=1.6em]
  \item \textbf{Current reads \texttt{0} or blows a fuse:} you used the \texttt{V$\Omega$mA}
        jack for current. Use the \textbf{10\,A} jack and the \textbf{10A} dial position.
  \item \textbf{Voltage shows \texttt{OL}:} wrong (too-low) range --- you're on DC
        \textbf{200}, which is correct for $\sim$22\,V; make sure it's DC, not AC.
  \item \textbf{Leading ``$-$'' sign:} leads are reversed. Fine --- swap them or just note
        the value.
  \item \textbf{Readings jump around:} passing clouds or the panel is warming/shading.
        Wait for steady sun.
  \item \textbf{Meter turned off by itself:} auto power-off after 15\,min --- just turn it
        back on.
\end{itemize}

\vspace{4pt}
\begin{center}\small Meter details: \texttt{datasheets/ABMUAM33DRD-EN-v1.pdf} (AstroAI AM33D).\end{center}

\end{document}
